\cotefile{157-2}
\scanpage{195}\mspage{140}
2) L'hom
\[
 \beta:\mathrm{id}_{\mathrm{HOT}_{\unsure{W|w}}(S)}\longleftarrow f_!f^*
\]
On peut décrire $\mathrm{HOT}_W(S)$ par $(\underline{\mathrm{Fib}}\,S)(W_S^d)^{-1}$, \inserted{ou $(\underline{\mathrm{Liss}}\,S)(W_S^d)^{-1}$} par \unsure{exemple}, de sorte que $f^*$ a une description simple par
\[
 f^*(\deleted{\unsure{h}}\mathrm{hot}_{S,W}(X))=\mathrm{hot}_{S',W}(X\times_S S')
\]
(pour $X$ lisse sur $S$). On a ensuite
\[
 f_! f^*(\mathrm{hot}_{S,W}(\unsure{x|X}))=\mathrm{hot}_{S,W}(X\times_S S'),
\]
mais attention, $X\times_S S'$ n'est pas a priori dans $\underline{\mathrm{Liss}}\,S$, seulement dans $\mathrm{Cat}_{/S}$ sans plus. Mais \unsure{cela} n'empêche qu'il donne ce qu'il faut, par localisation $(W_S^d)^{-1}$ de $\mathrm{Cat}_{/S}$. Pour définir $\beta_X$, il suffit donc de trouver une flèche
\[
 \beta_X:\unsure{X'|X}\longleftarrow X\times_S S'
\]
dans $\mathrm{Cat}_{/S}$, et on voit bien
% slanted
\margin{Quand $\xi\in\mathrm{HOT}_W(S)$ est défini par $X$ dans $\mathrm{Cat}_{/S}$ non néc. lisse, on \deleted{\unsure{prend}} doit trouver $\beta_X\unsure{{:}|{}}\Phi_S(X)\times_S S'\longrightarrow X$ dans $\mathrm{Cat}_{/S}(W_S^d)^{-1}$, on prend $(\unsure{i|1}_X)^{-1}\unsure{\mathrm{pr}_1^{\Phi_S(X)}|\mathrm{pr}_1^{\Phi_s(X)}}$ (cf. \unsure{page|pages} 143}
\[
 \left\{\begin{aligned}
 \alpha &: \mathrm{id}\longrightarrow f^*f_!\\
 \beta &: \mathrm{id}\longleftarrow f_!f^*
 \end{aligned}\right.
 \quad\mathrm{pr}_1^{\deleted{\unsure{X}}}.\qquad{}^{\unsure{+|*}}
\]

3) La compatibilité
\diagram{\begin{tikzcd}[ampersand replacement=\&]
 f_! \arrow[r,"f_!\star\alpha"] \arrow[rr,bend right=16,swap,"\mathrm{id}"] \& \deleted{\unsure{f^*}} f_! f^* f_! \arrow[r,"\beta\star f_!"] \& f_!
\end{tikzcd}}
On doit donc prendre $X'$ sur $S'$, et \deleted{\unsure{calculer}} expliciter les flèches à partir de l'objet
\[
 \xi'=\mathrm{hot}_{S',W}(X')
\]
% archive sheet number: 194
\scanpage{196}\mspage{141}
On a
\[
 f_!(\xi')=\mathrm{hot}_{S,\unsure{W|w}}(X'\deleted{\unsure{/S}})\deleted{\unsure{=}}\unsure{,|;}
\]
et pour calculer $f^*f_!(\xi')$, il faut comme \unsure{tantôt} prendre $\Phi_S(X')$ et on trouve
\[
 f^{\unsure{d|*}}f_!(\xi')=\deleted{\unsure{\mathrm{hot}}}\inserted{\mathrm{hot}_{S',W}}(\Phi_S(X')\times_S S').
\]
Enfin,
\[
 f_!f^*f_!(\xi')=\mathrm{hot}_{S,W}(\Phi_S(X')\times_S S')
\]
Il faut expliciter les flèches
\diagram{\begin{tikzcd}[ampersand replacement=\&,column sep=large]
 (*) \& \underset{\rule{.35pt}{.4ex}\rule{2em}{.35pt}\rule{.35pt}{.4ex}}{f_!(\xi')} \arrow[r,"f_!(\alpha_{\xi'})"] \& \underset{\rule{.35pt}{.4ex}\rule{3.35em}{.35pt}\rule{.35pt}{.4ex}}{f_!f^df_!(\xi')} \arrow[r,"\beta_{f_!(\xi')}"] \& \underset{\rule{.35pt}{.4ex}\rule{2em}{.35pt}\rule{.35pt}{.4ex}}{f_!(\xi')}\\
 \& X'/S \arrow[r,"\alpha_{X'}","{\substack{\text{déduit de}\\i_{X'}}}"'] \& \underset{\text{sur }S}{\Phi_S(X')\times_S S'} \arrow[r,dashed,"\beta_{X'}"] \& X'/S
\end{tikzcd}}
$\alpha_{X'}$ est déduit de $i_{X'}:X'\longrightarrow\Phi_S(X')$ comme l'unique $S'$-morphisme rendant commutatif le diagramme
\diagram{\begin{tikzcd}[ampersand replacement=\&]
 \& X' \arrow[r,"\alpha_{X'}"] \arrow[rr,bend left=15,"i_{X'}"] \arrow[d] \& \Phi_S(X')\times_S S' \arrow[r,"\mathrm{pr}_1"] \arrow[dl,dashed] \& \Phi_S(X') \arrow[d]\\
 (**) \& S' \arrow[rr,no head,very near end,"{\unsure{=}}"] \& \& S
\end{tikzcd}}
Pour décrire $\beta_{f_!(\xi')}$, où $f_!(\xi')=\mathrm{hot}_{S,W}(X')$, il faut donner une description \inserted{générale} \unsure{des|de} $\beta_\xi$ qui soit valable pour \uline{tout} représentant $X$ de $\xi$, pas seulement dans le cas $X$ lisse sur $S$. \deleted{\unsure{Donc}} Pour calculer $f^*(\xi)$ il faut donc prendre le $\mathrm{hot}_{S',W}$ de $\Phi_S(X)\times_S S'$ \deleted{\unsure{et puis}}, puis pour $f_!f^{\unsure{*}}(\xi)$, il
% archive sheet number: 195
\scanpage{197}\mspage{142}
faut prendre le $\mathrm{hot}_{S,\unsure{W|w}}$ du même $\Phi_S(X)\times_S S'$. \deleted{\unsure{Il}} Il faut définir, si on peut, une flèche
\[
 \Phi_S(X)\times_S S'\longrightarrow X\qquad\text{dans }\mathrm{Cat}_{/S}
\]
ou du moins dans $(\mathrm{Cat}_{/S})(W_S^d)^{-1}$. \deleted{\unsure{Mais}} Mais on a
\[
 i_X:X\longrightarrow\Phi_S(\unsure{x|X}),
\]
qui est dans $W_S^d$, et qui par produit cartésien définit
\[
 i_{\unsure{\alpha|X}}\times_S\mathrm{id}_{S'}:X\times_S S'\longrightarrow\Phi_S(X)\times_S S'\unsure{*|.|?}
\]
(en d'autres termes on a pris l'image inverse de $i_X$ par $S'\longrightarrow S$). \deleted{Mais \unsure{on a}} Cette \deleted{\unsure{vu que} le foncteur image inverse transforme $W_S^d$ en $W_{S'}^d$\unsure{)}} flèche \uline{est-elle \inserted{aussi} dans $W_S^d$} ?
% slanted
\margin{Non, elle ne l'est pas, car si $S'=\{s\}$, cela signifierait \unsure{(}que $X'_s\longrightarrow s\backslash X$ est dans $W$, ce qui est faux en général.}
Si oui, on a le diagramme
\diagram{\begin{tikzcd}[ampersand replacement=\&,column sep=huge]
 \Phi_S(X)\times_S S' \arrow[r,dashed,"\beta_X\,?"] \& X\\
 X\times_S S' \arrow[u,"i_X\times_S\mathrm{id}_{S'}","{=\pi_X}"'] \arrow[ur,"\mathrm{pr}_1^X"'] \&
\end{tikzcd}}
et on définit la flèche \inserted{cherchée} $\beta_X$ \uline{dans $(\mathrm{Cat}_{/S})(W_S^d)^{-1}$} par la commutativité de ce triangle, donc comme
\[
 \beta_X=\mathrm{pr}_1^X\circ\underset{\substack{\rule{.35pt}{.4ex}\rule{1.6em}{.35pt}\rule{.35pt}{.4ex}\\\text{inverse dans}\\(\mathrm{Cat}_{/S})(W_S^d)^{-1}}}{\pi_X^{-1}}
\]
% archive sheet number: 196
\scanpage{198}\mspage{143}
Mais ça ne marche pas -- cf \unsure{annotation} ! On va faire donc l'inverse
% slanted
\margin{\unsure{Description} de $\beta_\xi$ pour $\xi=\mathrm{hot}_{S,\unsure{W|w}}(X)$, avec $X$ quelconque dans $\mathrm{Cat}_{/S}$, à expliciter dans rédaction en forme.}
\diagram{\begin{tikzcd}[ampersand replacement=\&,column sep=huge]
 \Phi_S(X)\times_S S' \arrow[r,"{\mathrm{pr}_1^{\Phi_S(\unsure{X|Y})}}" ] \arrow[dr,dashed,"\beta_X\,?"'] \& \Phi_S(X)\\
 \& X \arrow[u,"{i_X\quad\text{dans }W_S^d}"']
\end{tikzcd}}
on définit $\beta_X$ comme $(\unsure{v|i}_X)^{-1}\mathrm{pr}_1^{\Phi_S(X)}$, où $(i_X)^{-1}$ est défini dans $(\mathrm{Cat}_{/S})(W_S^d)^{-1}$.

Pour prouver la commutativité de $(*)$ (p. 141), comme $v_{X'}:X'\longrightarrow\Phi_S(X')$ est dans $W_S^{\unsure{1|d}}$, il suffit de prouver, au lieu de \deleted{\unsure{$\mathrm{pr}_1^{\Phi_S(X')}\alpha_{X'}=\mathrm{id}_{X'}$}}, \inserted{$\beta_{X'}\alpha_{X'}=\mathrm{id}_{X'}$}, la relation \unsure{déduite} en composant avec $v_{X'}:X'\longrightarrow\Phi_S(X')$
\[
 \underset{\substack{\rule{.35pt}{.4ex}\rule{2.5em}{.35pt}\rule{.35pt}{.4ex}\\\mathrm{pr}_1^{\Phi_S(X')}}}{i_{X'}\beta_{X'}}\alpha_{X'}=i_{X'},
\]
i.e.
\[
 \mathrm{pr}_1^{\Phi_S(X')}\circ\alpha_{X'}=i_{X'},
\]
ce qui n'est autre que la relation de définition de $\alpha_{X'}$ (cf $(**)$ p. 141), on gagne !

4) Compatibilité
\diagram{\begin{tikzcd}[ampersand replacement=\&]
 f^* \arrow[r,"{\alpha\unsure{+|\star} f^*}"] \arrow[rr,bend right=16,"\mathrm{id}"'] \& f^*f_!f^* \arrow[r,"f^*\star\beta"] \& f^*,
\end{tikzcd}}
% archive sheet number: 197
\scanpage{199}\mspage{144}
Partons donc d'un\unsure{e}\deleted{\unsure{e}} $\xi\in\mathrm{Ob}\,\mathrm{HOT}_{\unsure{W|w}}(S)$, qu'on peut supposer défini par $X\in\mathrm{Ob}\,\underline{\mathrm{Fib}}\,S$. Alors
\[
 \begin{aligned}
 \deleted{\unsure{f^*}}\qquad f^*(\xi)&=\deleted{\unsure{\mathrm{Hot}}}\mathrm{hot}_{S',W}(X\times_S S'),\\
 f_!f^*(\xi)&=\mathrm{hot}_{S,W}(X\times_S S')\\
 f^*f_!f^*(\xi)&=\mathrm{hot}_{S',W}(\Phi_S(X\times_S S')\times_S S').
 \end{aligned}
\]
Il faut expliciter maintenant
\diagram{\begin{tikzcd}[ampersand replacement=\&,column sep=small]
 \underset{\substack{\unsure{\Vert|''}\\\mathrm{hot}_{S',W}(X')}}{f^*(\xi)} \arrow[r,"\alpha_{f^*(\xi)}"] \& \underset{\mathrm{hot}_{S',W}(\Phi_S(X')\times_S S')}{\underset{\rule{.35pt}{.4ex}\rule{4em}{.35pt}\rule{.35pt}{.4ex}}{f^{\unsure{d|*}}f_!f^*(\xi)}} \arrow[r,"f^*(\beta_\xi)"] \& \underset{\mathrm{hot}_{S',W}(X')}{\underset{\rule{.35pt}{.4ex}\rule{2.1em}{.35pt}\rule{.35pt}{.4ex}}{f^*(\xi)}},
\end{tikzcd}}
et vérifier que le composé est l'identité. \deleted{\unsure{Comme}} On a \unsure{posé|pris}
\[
 X'=X\times_S S',
\]
\deleted{\unsure{Or, l'on aurait}} et \deleted{\unsure{tout}} le diagramme est donc aussi
\diagram{\begin{tikzcd}[ampersand replacement=\&,column sep=large]
 X' \arrow[r,"{\alpha_{X'\deleted{\unsure{\xi}}}}"] \arrow[rr,bend left=12,"\mathrm{id}_{X'}"] \& \Phi_S(X')\times_S S' \arrow[r,"f^*(\beta_\xi)"] \& X'
\end{tikzcd}}
dans $(\mathrm{Cat}_{/S'})(W_{S'}^d)^{-1}$. Pour vérifier la commutativité, \unsure{on peut} composer avec
\[
 X\times_S S'=X'\xrightarrow{i_X\times_S\mathrm{id}_{S'}}\Phi_S(X)\times_S S'
 \qquad(\text{où }i_X:X\longrightarrow\Phi_S(X))
\]
% slanted
\margin{\uline{NB} Comme $X$ est lisse $/S$, et $\Phi_S(X)$ \unsure{a fortiori}, il s'ensuit que $X\longrightarrow\Phi_S(X)$ \unsure{reste} une $W$-équiv. colocale par tt \unsure{chgt} de base}
\unsure{qui} est dans $W_{S'}^d$, or on vérifie, sur la construction générale de $\beta_\xi=\beta_X$ (page précédente), que l'on a comm. du
\diagram{\begin{tikzcd}[ampersand replacement=\&,column sep=huge]
 \Phi_S(X')\times_S S' \arrow[r,"{f^*(\deleted{\unsure{\beta}}\beta_\xi)}"] \arrow[dr,"{\Phi_S(\mathrm{pr}_1^X)\times_S S'}"'] \& X' \arrow[d,"{i_X\times_S S'}"]\\
 \& \Phi_S(X)\times_S S'
\end{tikzcd}}
% archive sheet number: 198
\scanpage{200}\mspage{145}
On est donc ramené à la commutativité de
\diagram{\begin{tikzcd}[ampersand replacement=\&,column sep=large]
 \& X\times_S S'=X' \arrow[r,"\alpha_{X'}"] \arrow[dr,"{i_X\times_S S'}"'] \& \Phi_S(X')\times_S S' \arrow[d,"{\Phi_S(\mathrm{pr}_1^X)\times_S S'}"]\\
 (\unsure{*|w}) \& \& \Phi_S(X)\times_S S' \arrow[ul,phantom,very near start,"?"]
\end{tikzcd}}
dans $(\mathrm{Cat}_{/S'})(W_{S'}^d)^{-1}$.

\deleted{Or la flèche \inserted{\unsure{de|dans} $\mathrm{Cat}_{/S}$} \deleted{$\unsure{\Phi_S(\mathrm{pr}_1^X)\times_S S'}$} \inserted{$\Phi_S(\mathrm{pr}_1^X):\underset{\substack{\rule{.35pt}{.4ex}\rule{2em}{.35pt}\rule{.35pt}{.4ex}\\Z'}}{\Phi_S(X')}\xrightarrow[u]{}\underset{\substack{\rule{.35pt}{.4ex}\rule{1.8em}{.35pt}\rule{.35pt}{.4ex}\\Z}}{\Phi_S(X)}$} donne, par fonctorialité de $Z\longmapsto(Z\times_S S'\xrightarrow{\mathrm{pr}_{\unsure{11|1}}^Z}Z)$ de $\mathrm{Cat}_{/S}$ dans $\underline{\mathrm{Fl}}(\underline{\mathrm{Cat}}_{/S})$, le carré commutatif
\diagram{\begin{tikzcd}[ampersand replacement=\&]
 \deleted{\unsure{Z}} Z'\times_S S' \arrow[r,"\mathrm{pr}_1^{Z'}"] \arrow[d,"{u\times_S S'}"'] \& Z' \arrow[d,"u"]\\
 Z\times_S S' \arrow[r,"\mathrm{pr}_1^Z"] \& Z
\end{tikzcd}}
i.e.
\diagram{\begin{tikzcd}[ampersand replacement=\&,column sep=large]
 \Phi_S(X')\times_S S' \arrow[r,"{\mathrm{pr}_1^{\Phi_S(X')}}"] \arrow[d,"{\Phi_S(\mathrm{pr}_1^X)\times_S S'}"'] \& \Phi_S(X') \arrow[d,"{\Phi_S(\mathrm{pr}_1^X)}"]\\
 \Phi_S(X)\times_S S' \arrow[r,no head,"{\mathrm{pr}_1^{\Phi_S(X)}}"] \& \Phi_S(X)
\end{tikzcd}}}

Essayons de prouver que c'est comm. dans $(\mathrm{Cat}_{/S'})\deleted{\unsure{\times_S}}$, pas seulement dans \unsure{le|la} localisé par $W_{S'}^d$. On peut le regarder comme diagramme dans $\mathrm{Cat}_{/S}$ et prouver la comm. à ce titre. Cette relation de \unsure{commutativité} $vu=w$ se décompose alors en deux, suivant les deux facteurs $\Phi_S(X)$ et $S'$ de l'objet but $\Phi_S(X)\times_S S'$. Suivant le facteur $S'$, la relation signifie simplement que $\alpha_{X'}$ est un $S'$-morphisme, c'est OK. Pour le facteur $\Phi_S(X)$, on \deleted{\unsure{utilise}} prolonge le diagramme $(*)$ à droite par le carré
\diagram{\begin{tikzcd}[ampersand replacement=\&,column sep=large]
 X' \arrow[r,dashed,"\alpha_{X'}"] \arrow[dr,dashed] \& \Phi_S(X')\times_S S' \arrow[r,"{\mathrm{pr}_1^{\Phi_S(X')}}"] \arrow[d,"{\Phi_S(\mathrm{pr}_1^X)\times_S S'}"'] \& \Phi_S(X') \arrow[d,"{\Phi_S(\mathrm{pr}_1^X)}"]\\
 \& \Phi_S(X)\times_S S' \arrow[r,"{\mathrm{pr}_1^{\Phi_S(X)}}"'] \& \Phi_S(X)
\end{tikzcd}}
% archive sheet number: 199
\scanpage{201}\mspage{146.}
et \unsure{compte} tenu que le composé de la première ligne horizontale est $i_{X'}$, par construction de $\alpha_{X'}$, on est ramené à la commutativité de
\diagram{\begin{tikzcd}[ampersand replacement=\&,column sep=large]
 \& X\times_S S'=X' \arrow[r,no head,"i_{X'}"] \arrow[d,"{\unsure{j|i}_X\times_S S'}"'] \& \Phi_S(X') \arrow[d,"{\Phi_S(\mathrm{pr}_1^X)}"]\\
 (*) \& \Phi_S(X)\times_S S' \arrow[r,"{\mathrm{pr}_1^{\Phi_S(X)}}"'] \& \Phi_S(X).
\end{tikzcd}}
Or le morphisme $i_X:X\longrightarrow\Phi_S(X)$ dans $\mathrm{Cat}_{/S}$ donne lieu au carré commutatif (comme \unsure{toute} flèche $i:X\longrightarrow Y$)
\diagram{\begin{tikzcd}[ampersand replacement=\&,column sep=large]
 X\times_S S' \arrow[r,"\mathrm{pr}_1^X"] \arrow[d,"{i_X\times_S S'}"'] \& X \arrow[d,"\unsure{i'_X|i_X}"]\\
 \Phi_S(X)\times_S S' \arrow[r,"{\mathrm{pr}_1^{\Phi_S(X)}}"] \& \Phi_S(X),
\end{tikzcd}}
de sorte que le composé inférieur de $(*)$ n'est autre que celui de $(**)$, donc égal au composé sup. de ce dernier, et la comm. de $(*)$ revient à celle de
\diagram{\begin{tikzcd}[ampersand replacement=\&,column sep=large]
 X\times_S S'=X' \arrow[r,"i_{X'}"] \arrow[d,"\mathrm{pr}_1^X"'] \& \Phi_S(X') \arrow[d,"{\Phi_S(\mathrm{pr}_1^X)}"]\\
 X \arrow[r,"i_X"] \& \Phi_S(X),
\end{tikzcd}}
qui est évidente par fonctorialité de \inserted{$\underbrace{(Z\xrightarrow{i_Z}\Phi_S(Z))}_{\deleted{\text{\unsure{l'hom}}}\,\text{\unsure{en }}Z}$}, appliquée à la flèche $\mathrm{pr}_1^X:X'\longrightarrow X$ de $\mathrm{Cat}_{/S}$.
% archive sheet number: 200
\scanpage{202}\mspage{14\deleted{\unsure{6}}7}
\deleted{\unsure{tout|pour} $Y$ \unsure{de|dans} $\mathrm{Cat}_{/S}$, $X\times_S Y\longrightarrow X'\times_S Y$ est dans $W_S^d$.}

On a ainsi prouvé que les foncteurs
\diagram{\begin{tikzcd}[ampersand replacement=\&,column sep=large]
 (*) \& \mathrm{HOT}_{\unsure{W|w}}(S') \arrow[r,shift left=2pt,"f_!"] \& \mathrm{HOT}_W(S) \arrow[l,shift left=2pt,"f^*"]
\end{tikzcd}}
sont adjoints l'un de l'autre, avec des morphismes d'adjonction
\[
 \left\{\begin{aligned}
 \mathrm{id}&\xrightarrow{\alpha} f^*f_!\\
 \mathrm{id}&\xleftarrow{\beta} f_!f^*
 \end{aligned}\right.
\]
précisés plus haut. Je présume que la même démonstration, en travaillant avec les \deleted{\unsure{relatifs}} localiseurs \inserted{grossiers} $W_{\unsure{S\to|S'|S}},W_{S'}$ au lieu de $W_S^d,W_{S'}^d$, donnerait que les foncteurs correspondants
\diagram{\begin{tikzcd}[ampersand replacement=\&,column sep=large]
 (**) \& \mathrm{HOT}_W^{\mathrm{lc}}(S') \arrow[r,shift left=2pt,"f_!^{\mathrm{lc}}"] \& \mathrm{HOT}_W^{\mathrm{lc}}(S) \arrow[l,shift left=2pt,"f^*_{\mathrm{lc}}"]
\end{tikzcd}}
sont adjoints l'un de l'autre, avec essentiellement ``les mêmes'' morphismes d'adjonction $\alpha,\beta$. Je n'ai pas envie de reprendre les raisonnements précédents, et préfère le \deleted{\unsure{prouver}} \uline{déduire} de l'adjonction dans la situation $(*)$. Cela me donnera l'occasion de mieux comprendre les traits formels essentiels de la situation. Je poserai pour simplifier
% archive sheet number: 201
\scanpage{203}\mspage{148}
\[
 \begin{aligned}
 \mathrm{HOT}_{\unsure{W|w}}(S)&=\unsure{A|\mathcal{A}},&\mathrm{HOT}_W(S')&=A'\\
 \mathrm{HOT}^{\mathrm{lc}}_W(S)&=A_0,&\mathrm{HOT}^{\mathrm{lc}}_W(S')&=A'_0,
 \end{aligned}
\]
et j'utilise les couples de foncteurs adjoints
\deleted{\diagram{\begin{tikzcd}[ampersand replacement=\&]
 A_0 \arrow[r] \& A
\end{tikzcd}}}
\diagram{\begin{tikzcd}[ampersand replacement=\&]
 A \arrow[r,shift left=2pt,"{\unsure{\rho|p}}"] \& A_0 \arrow[l,shift left=2pt,"i"] ,\qquad \& A' \arrow[r,shift left=2pt,"\rho'"] \& A'_0 \arrow[l,shift left=2pt,"i'"]
\end{tikzcd}}
où $\rho,\rho'$ sont des foncteurs \inserted{de} localisation, i.e. $i,i'$ des foncteurs pl. fidèles. \deleted{(\unsure{mais je ne} vais pas utiliser ce fait \inserted{de façon essentielle}).} Je dispose d'un couple de foncteurs adjoints
\diagram{\begin{tikzcd}[ampersand replacement=\&]
 A' \arrow[r,shift left=2pt,"f_!"] \& A \arrow[l,shift left=2pt,"f^*"],
\end{tikzcd}}
et je suppose \deleted{que} de plus

a) $\deleted{\unsure{f|f_*}} f_!$ est compatible avec les localiseurs de $A',A\unsure{{}^{*}|{}}$ définis par $\rho',\rho$ respectivement, donc définit un foncteur (bien déterminé à isom unique près)
\[
 f_!^{\unsure{0|\circ}}:A'_0\longrightarrow A_0,\qquad\deleted{\unsure{donc}}
\]
\deleted{\unsure{b) $f^*$}} donnant donc lieu à un diagramme commutatif (à isom \unsure{ponctuel|fonctoriel|fonctorel} près)
\diagram{\begin{tikzcd}[ampersand replacement=\&]
 \& A' \arrow[r,"\rho'"] \arrow[d,"f_!"'] \& A'_0 \arrow[d,"f_!^0"]\\
 (\mathrm{I}) \& A \arrow[r] \& A_0 .
\end{tikzcd}}
b) ``Dualement'', $f^*\inserted{{}:A\longrightarrow A'}$ est compatible avec les sous-catégories pleines \deleted{définies par} \inserted{images} essentielles des foncteurs pl. fid. $i,i'$, et induit
% archive sheet number: 202
\scanpage{204}\mspage{149}
\deleted{\unsure{donc}} \inserted{donc} un foncteur, \unsure{définit|défini} à isom unique près :
\[
 f_{\unsure{0|\circ}}^*:\unsure{A|\mathcal{A}}_0\longrightarrow A'_0,
\]
donnant lieu à un carré de foncteurs, commutatif à isom fonctoriel canonique près
\diagram{\begin{tikzcd}[ampersand replacement=\&]
 \& A_0 \arrow[r,"i"] \arrow[d,"f_0^*"'] \& A \arrow[d,"f^*"]\\
 (\mathrm{II}) \& A'_0 \arrow[r,"i'"] \& A'.
\end{tikzcd}}
Sous ces conditions, le couple de foncteurs
\diagram{\begin{tikzcd}[ampersand replacement=\&]
 A_0 \arrow[r,shift left=2pt,"f_!^0"] \& A'_0 \arrow[l,shift left=2pt,"f_0^*"]
\end{tikzcd}}
est un couple de foncteurs adjoints. (NB quand je parle de couples de foncteurs adjoints, celui écrit au dessus de l'autre est l'adjoint à gauche.)

\uline{Dém} Soient
\[
 \deleted{\unsure{\xi}}\unsure{\eta|y}_0\in\deleted{\unsure{\mathrm{Ob}}}\mathrm{Ob}\,A_0,\quad \deleted{\unsure{\eta}}\xi'_0\in\mathrm{Ob}\,A'_0,
\]
il faut définir une bijection bifonctorielle
\[
 \mathrm{Hom}_{A'_0}(f_!^0(\xi'_0),\eta_0)\overset{?}{\simeq}
 \mathrm{Hom}_{A^{\unsure{\deleted{*}|'}}_0}(\xi'_0,f_0^*(\eta_0)).
\]
Pour cela, on va écrire $\unsure{\xi_0|\xi'_0}$ sous la forme
\[
 \xi'_0=\unsure{\rho|p}'(\xi')\qquad \xi'\in\mathrm{Ob}\,A',
\]
il suffira d'établir la bijection bifonctorielle en remplaçant $\xi'_0$ par $\rho'(\xi')$, i.e.
\[
 \mathrm{Hom}_{A^{\deleted{\unsure{\prime}}}_0}(f_!^0(\rho'(\xi')),\eta_0)\simeq
 \mathrm{Hom}_{A'_0}(\rho'(\xi'),f_0^{\unsure{d|*}}(\eta_0)).
\]
% archive sheet number: 203
\scanpage{205}\mspage{150}
(En effet, la formule d'adjonction\deleted{\unsure{ici}}, s'interprète comme un isom. de bifoncteurs, ou ce qui revient au $=$ de foncteurs
\diagram{\begin{tikzcd}[ampersand replacement=\&]
 \unsure{A|\mathcal{A}}'_0 \arrow[r,shift left=2pt,"{\unsure{\varphi|\phi}}"] \arrow[r,shift right=2pt,"\psi"'] \& A_0^{\wedge},
\end{tikzcd}}
or
\[
 \begin{aligned}
 \underline{\mathrm{Hom}}(A'_0,A_0^{\wedge})&\longrightarrow\underline{\mathrm{Hom}}(A',A_0^{\wedge})\\
 \varphi&\longmapsto\varphi\circ\unsure{\rho|p}'
 \end{aligned}
\]
est pleinement fidèle, puisque $\rho'$ est un foncteur de localisation, donc il revient au $=$ de se donner un isom $\varphi\simeq\psi$, ou $\varphi\circ\rho'\simeq\psi\circ\rho'$). On doit donc transformer l'expression
\[
\begin{aligned}
 &\mathrm{Hom}_{A^{\deleted{\unsure{\prime}}}_0}(\underset{\substack{\rule{.35pt}{.4ex}\rule{4em}{.35pt}\rule{.35pt}{.4ex}\\\simeq\,\rho f_!(\xi')}}{f_!^{\unsure{0|\circ}}(\rho'(\xi')},\unsure{\eta|y}_0)
 \underset{\substack{\vert\\\text{comm.}\\\text{de (I)}}}{\simeq}\mathrm{Hom}_{A^{\deleted{\unsure{\prime}}}_0}(\rho f_!(\xi'),\eta_0)\\
 &\underset{\substack{\vert\\\text{adj. }\rho,\unsure{i}}}{\simeq}\mathrm{Hom}_{A_{\deleted{\unsure{0}}}^{\deleted{\unsure{\prime}}}}(f_!(\xi'),\deleted{\unsure{\rho^*}}i(\eta_0))
 \underset{\substack{\vert\\\text{adj. }f_!,f^*}}{\simeq}\mathrm{Hom}_{A'}(\xi',\underset{\substack{\rule{.35pt}{.4ex}\rule{2.8em}{.35pt}\rule{.35pt}{.4ex}\\i'(f_0^*\eta_0)}}{f^*i(\eta_0)})\\
 &\underset{\substack{\vert\\\text{comm. de}\\\text{(II)}}}{\simeq}\mathrm{Hom}_{A'}(\xi',i'(f_0^*(\eta_0)))
 \underset{\substack{\vert\\\text{adj }\rho',i'}}{\simeq}\mathrm{Hom}_{A'_0}(\underset{\substack{\rule{.35pt}{.4ex}\rule{2em}{.35pt}\rule{.35pt}{.4ex}\\\xi'_0}}{\rho'(\xi')},f_0^{\unsure{\alpha|d|*}}(\eta_0)),
\end{aligned}
\]
\deleted{\unsure{ce}} ce qui est la bijection annoncée.

Je résume l'essentiel des résultats obtenus sur $f_!,f_!^{\mathrm{lc}}$ :

\uline{Théorème 2} Soit $f:S'\longrightarrow S$ une flèche \unsure{dans|de} $\mathrm{Cat}$, considérons les foncteurs
\diagram{\begin{tikzcd}[ampersand replacement=\&,column sep=large]
 \deleted{\unsure{f}}\,\mathrm{HOT}_{\unsure{W|w}}(S') \arrow[r,shift left=2pt,"f_!"] \& \mathrm{HOT}_W(S) \arrow[l,shift left=2pt,"f^*"]\\
 \mathrm{HOT}_W^{\mathrm{lc}}(S') \arrow[r,shift left=2pt,"f_!^{\mathrm{lc}}"] \& \mathrm{HOT}_W^{\mathrm{lc}}(S) \arrow[l,shift left=2pt,"f_{\mathrm{lc}}^*"],
\end{tikzcd}}
% archive sheet number: 204
\scanpage{206}\mspage{151}
où les foncteurs $f_!,f_!^{\mathrm{lc}}$ sont décrits dans le th. 1 (p. 130). Les morphismes d'adjonction
\[
 \begin{aligned}
 \mathrm{id}&\longrightarrow f^*f_!\\
 \mathrm{id}&\longleftarrow f_!f^*
 \end{aligned}
\]
sont décrits p. 138 à 140.

On voudrait expliciter la formule \unsure{sympa}, obtenue en passant
\[
 \boxed{f_!f^*(\xi)=\mathrm{hot}_{S,\unsure{W|w}}(S'/S)\times_S\xi},
\]
pour ceci il faut expliciter le deuxième membre. Notons la

\uline{Proposition 4} Il existe un unique foncteur
\[
 (\xi,\unsure{\eta|y})\longmapsto\xi\times_S\eta:\mathrm{HOT}_W(S)\times\mathrm{HOT}_W(\deleted{\unsure{(|S}}S)\longrightarrow\mathrm{HOT}_W(S),
\]
% slanted
\margin{\uline{NB} La multiplication $\times_S$ dans $\mathrm{HOT}_W(S)$ \unsure{est} associative, commutative, unitaire etc etc, puisque c'est le cas dans $\underline{\mathrm{Liss}}\,S$.}
avec la propriété que pour $X,Y$ variables dans $\underline{\mathrm{Liss}}\,S$, on ait un isomorphisme bifonctoriel en $X,Y$
\[
 \mathrm{hot}_{S,W}(X)\times_S\mathrm{hot}_{S,W}(\unsure{\mathrm{hot}|w|{}}(\deleted{\unsure{x}}Y))\simeq\mathrm{hot}_{S,W}(X\times_S Y).
 \tag*{$(+)$}
\]
C'est formel, en utilisant seulement que sur $\underline{\mathrm{Liss}}\,S$, le bifoncteur
\[
 (X,Y)\longmapsto X\times_S Y
\]
est compatible avec les localiseurs $W_{\unsure{C|S}}^d\times W_S^d$ et $W_S^d$. Pour le voir, on est ramené à voir que si $X\xrightarrow{u}X'$ dans $\underline{\mathrm{Liss}}\,S$ est dans $W_S^d$, alors il en est de $=$ de $X\times_S Y\xrightarrow{u\times_S\mathrm{id}_Y}X'\times_S Y$, \deleted{\unsure{En}} pour tt
% archive sheet number: 205
\scanpage{207}\mspage{152}
$Y$ de $\underline{\mathrm{Liss}}\,S$. En fait, on n'utilise que la lissité de $X,X'$, pas celle de $Y$, et c'est immédiat, car \deleted{d'après le} si $X\longrightarrow X'$ est dans $W_S^d$, $X\times_S Y\longrightarrow X'\times_S Y^{\deleted{\unsure{S}}}$ est dans $W_Y^d$, \deleted{\unsure{cf. th. 1}\ \illegible{2}} donc aussi dans $W_S^d$ (cf th. 1 appliqué à $Y\longrightarrow S$). Cette remarque \deleted{\unsure{prouve}} \inserted{suggère} en même temps :

\uline{Corollaire 1} \inserted{(\uuline{Faux}\unsure{!})} La formule bifonctorielle $(\unsure{+|*})$ est valable pour $X$ \deleted{lisse,} dans $\underline{\mathrm{Liss}}\,S$, $Y$ quelconque dans $\mathrm{Cat}_{/S}$, et aussi dans le cas inverse.
% slanted
\margin{Cet énoncé est sûrement \uuline{faux} ! \unsure{Contrexemple|Contreexemple} p. 153. Cf cor. 1 rectifié page \unsure{154|155}\deleted{\unsure{155}}}

Dans le premier cas, considérons
\[
 Y\xrightarrow{i_Y}\Phi(Y),\quad\text{dans }W_S^d\text{ donc}
\]
\deleted{\unsure{aussi}} iso dans $(\mathrm{Cat}_{/S})(W_S^d)^{-1}$, on aura donc
\[
\begin{aligned}
 \mathrm{hot}_{S,\unsure{W|w}}(X)\times_S\mathrm{hot}_{S,W}(Y)&\xrightarrow{\sim}\mathrm{hot}_{S,W}(X)\times_S\mathrm{hot}_{S,W}(\Phi(Y))\\
 &\overset{\substack{\text{car }X\ \text{\uline{et}}\ \Phi(Y)\ \text{lisses sur }S}}{\xrightarrow{\sim}}\mathrm{hot}_{S,W}(X\times_S\Phi(Y)),
\end{aligned}
\]
\deleted{\unsure{or}} je dis que
\[
 X\times_S Y\xrightarrow{X\times_S i_Y}X\times_S\Phi(Y)
\]
est encore dans $W_S^d$. Il suffit de voir que c'est dans $W_X^d$ (cf th. 1), \deleted{\unsure{donc}} \inserted{et pour ceci,} que le chgt. de base lisse $X\longrightarrow S$ induit $\mathrm{Cat}_{/S}\longrightarrow\mathrm{Cat}_{/\unsure{x|X}}$ compatible avec les localiseurs $W_S^d,W_X^d$. \deleted{Ce point étant admis, \unsure{donc}\ \illegible{1}} \inserted{Mais c'est \uline{faux} (prendre $X=\{s\}$) :}
% archive sheet number: 206
\scanpage{208}\mspage{153}
\deleted{\[
 \mathrm{hot}_{S,\unsure{W|w}}(X\times_S\Phi(Y))\xleftarrow{\sim}\mathrm{hot}_{S,W}(X\times_S Y),
\]
d'où la formule annoncée.

Il reste à prouver séparément le résultat suivant, qui m'avait échappé quand j'avais examiné les propriétés générales des foncteurs $f^{\unsure{d|*}}$ pour les $\mathrm{HOT}_W$ :

\uline{Proposition} Soit $f:S'\longrightarrow S$ un morphisme \uline{lisse} dans $\mathrm{Cat}$. Alors
\[
 f^d:X\longmapsto X\times_S S':\mathrm{Cat}_{/S}\longrightarrow\mathrm{Cat}_{/S'}
\]
\unsure{est} compatible avec les localiseurs $W_S^d,W_{S'}^d$.}

Montrons que le corollaire est faux. Je vais prendre $Y=\{s\}$, catégorie ponctuelle en $s$, \deleted{et \unsure{pour}} alors $\Phi(Y)$ est la cat. fibrée à fibres discrètes définie par le foncteur représentable \inserted{$F$} défini par $s$, donc c'est $S_{/F}$ et, en l'occurrence $S_{/s}$ :
\[
 \Phi(\{s\})\simeq S_{/s}\qquad\text{comme catégorie sur }S.
\]
\unsure{Si} p. ex. $s$ est élément final \inserted{$e$} de $S$, on trouve
\[
 \Phi(e_S)\simeq S
\]
donc
\[
 \mathrm{hot}_{S,W}(X)\times_S\mathrm{hot}_{S,W}(e)=\mathrm{hot}_{S,W}(X)
\]
d'autre part
\[
 \mathrm{hot}_{S,W}(X\times_S e)=\mathrm{hot}_{S,W}(X_e).
\]
% archive sheet number: 207
\scanpage{209}\mspage{154}
Si les deux étaient \deleted{\unsure{égaux}} \inserted{iso dans $\mathrm{HOT}_{\deleted{\unsure{S}}\unsure{W|w}}(S)$} par la flèche canonique
\[
 X_e\longrightarrow X,
\]
en appliquant le foncteur $p_!$ (où $p:S\longrightarrow e$) il s'ensuivrait que $X_e\longrightarrow X$ est dans $W$. Si par exemple $S=\unsure{\Delta|\mathbb{\Delta}}^1$, cela impliquerait que pour tout $X$ lisse sur $S$, $X_1\longrightarrow X$ est \deleted{dans} $W$, alors que visiblement la lissité implique que $X_0\longrightarrow X$ est dans $W$. \deleted{Si la lissité implique} Donc ça reviendrait à dire que \deleted{$X$ lisse} $\underline{\mathrm{Liss}}\,\Delta^1=\deleted{\unsure{\mathrm{Fib}\Delta^1}}\underline{\mathrm{Liss}}_0\unsure{\{\Delta^1\}|\Delta^1}$, d'où aussitôt $\underline{\mathrm{Liss}}\,S=\underline{\mathrm{Liss}}_0\,S$ \deleted{\unsure{ça se}} pour \deleted{tout} $S$ dans $\mathrm{Cat}$, ça se saurait ! Le contre \unsure{exemple} le plus trivial est de prendre $X=\{0\}\lhook\joinrel\longrightarrow S=\Delta^1$, l'inclusion de l'élément ``générique'', immersion ouverte, donc lisse, la fibre \deleted{\unsure{en}} $X_1$ est vide !

\margin{$\unsure{\mathrm{Z}|2}$}
\uline{Corollaire 1 rectifié} Si $X$ est dans $\mathrm{Cat}_{/S}$, pas néc. lisse sur $S$, alors le foncteur
\[
 \mathrm{mult}_{X/S}:Y\unsure{\longrightarrow|\longmapsto} Y\times_S X:\deleted{\unsure{\mathrm{Cat}}}\mathrm{Cat}_{/S}\longrightarrow\mathrm{Cat}_{/S}
\]
induit un foncteur
\[
 \underline{\mathrm{Liss}}\,S\longrightarrow\mathrm{Cat}_{/S}
\]
compatible avec les localiseurs $W_S^d,W_S^d$, donc induit un foncteur
\[
 (\underline{\mathrm{Liss}}\,S)(W_S^d)^{-1}\longrightarrow(\mathrm{Cat}_{/S})(W_S^d)^{-1}
\]
% archive sheet number: 208
\scanpage{210}\mspage{155}
ou encore \unsure{(compte|compte} tenu de l'équivalence \inserted{induite par l'inclusion}
\[
 (\underline{\mathrm{Liss}}\,S)(W_S^d)^{-1}\xrightarrow{\approx}(\underline{\underline{\mathrm{Cat}}}_{/S})(W_S^d)^{-1}
\]
un foncteur, encore noté $\mathrm{mult}^{\mathrm{hot}}_{\deleted{\unsure{X/S}}X/S}$
\[
 \mathrm{mult}^{\mathrm{hot}}_{X/S}:\mathrm{HOT}_{\unsure{W|w}}(S)\unsure{-|\longrightarrow}\mathrm{HOT}_W(S).
\]
\deleted{Sa valeur} Ce foncteur satisfait donc la relation
\[
 \mathrm{mult}^{\mathrm{hot}}_{X/S}(\mathrm{hot}_{S,W}(Y))\overset{\mathrm{iso.can}}{\simeq}\mathrm{hot}_{S,W}(X\times_S Y)
\]
pour $Y$ \uline{lisse sur} $S$,
\margin{$\unsure{\mathrm{Z}|2}$}
et pour $Y$ quelconque dans $\mathrm{Cat}_{/S}$, \unsure{à|''}
\[
 \mathrm{mult}^{(\mathrm{hot})}_{X/S}(\mathrm{hot}_{S,W}(Y))\overset{\mathrm{iso\,can}}{\simeq}\mathrm{hot}_{S,W}(X\times_S\Phi_S(Y)).
\]
\deleted{Mais on se gardera de confon\unsure{dre}} Quand $X$ est lui-même lisse sur $S$, ce foncteur \deleted{\unsure{est donc}} n'est autre que la multiplication par $\mathrm{hot}_{S,W}(X)$ :
\[
 \mathrm{mult}^{(\mathrm{hot})}_{X/S}=\mathrm{mult}_{\mathrm{hot}_{S,W}(X)}
 \qquad\text{si }X\text{ \uline{lisse sur }}S.
\]
\inserted{Mais} \unsure{Sa|sa} valeur pour \deleted{$\unsure{Y}$ \unsure{lisse sur $S$}} $\mathrm{hot}_{S,W}(Y)$\unsure{,|$\ast$} quand $Y$ n'est \uline{pas} lisse sur $S$, ne doit pas être confondue avec $\mathrm{hot}_{S,W}(X\times_S Y)$.
\margin{$\mathrm{Z}$}

Ainsi, si $X$ et $Y$ sont dans $\mathrm{Cat}_{/S}$, considérons :
\[
 X\times_S Y\qquad\text{et}\qquad
 \tikz[baseline=(h.base)]{\node[ellipse,inner xsep=4pt,inner ysep=2pt] (h) {$\mathrm{hot}_{S,W}(X\times_S Y).$}; \draw (h.18) to[out=155,in=5] (h.north) to[out=180,in=90] (h.west) to[out=-90,in=180] (h.south) to[out=0,in=-90] (h.east) to[out=90,in=0] (h.6);}
\]
a) Si ni $X$, ni $Y$ ne sont supposés lisses sur $S$, \unsure{ce dernier|ça} ``\unsure{n'a} aucune interprétation \unsure{naturelle} en termes d'opérations \deleted{entre catégories} \inserted{sur \deleted{des} catégories} \deleted{dans des} $\mathrm{HOT}_W(\unsure{S|s})$''.\textsuperscript{\unsure{$\ast$}}
\margin{\unsure{Sauf} dans le cas où l'une des deux est dans $\underline{\unsure{\mathrm{Fibhot}|\mathrm{FibHot}}}\,S$, \unsure{comme on vient de voir}.}
% slanted
% archive sheet number: 209
